\documentclass[10pt,a4paper]{amsart}
\usepackage{amsmath,amssymb,amsthm,graphicx}
% Retained minimal preamble for the Stacks Project chapter more-algebra.tex (More on Algebra).
% The chapter file pinned-original-source.tex carries only \input{preamble} on its first line, so no
% preamble text lives inside the chapter itself.  The Stacks Project uses one single shared preamble for
% every chapter; the snapshot stages it identically as sources/stacks-chow/preamble.tex and
% sources/stacks-etale-cohomology/preamble.tex.  The block below is copied mechanically, line for line,
% from sources/stacks-chow/preamble.tex lines 14-18 (the Xy-pic setup) and 154-217 (the native theorem
% declarations together with the mathematical macros).  Nothing was added, removed or redefined; the
% wrapper class and the amsmath/amssymb/amsthm/graphicx base packages are supplied by the assembler.
% Chapter: more-algebra.tex -- More on Algebra.
% For commutative diagrams we use Xy-pic
\usepackage[all]{xy}

% We use 2cell for 2-commutative diagrams.
\xyoption{2cell}

% Theorem environments.
%
\theoremstyle{plain}
\newtheorem{theorem}[subsection]{Theorem}
\newtheorem{proposition}[subsection]{Proposition}
\newtheorem{lemma}[subsection]{Lemma}

\theoremstyle{definition}
\newtheorem{definition}[subsection]{Definition}
\newtheorem{example}[subsection]{Example}
\newtheorem{exercise}[subsection]{Exercise}
\newtheorem{situation}[subsection]{Situation}

\theoremstyle{remark}
\newtheorem{remark}[subsection]{Remark}
\newtheorem{remarks}[subsection]{Remarks}

\numberwithin{equation}{subsection}

% Macros
%
\def\lim{\mathop{\mathrm{lim}}\nolimits}
\def\colim{\mathop{\mathrm{colim}}\nolimits}
\def\Spec{\mathop{\mathrm{Spec}}}
\def\Hom{\mathop{\mathrm{Hom}}\nolimits}
\def\Ext{\mathop{\mathrm{Ext}}\nolimits}
\def\SheafHom{\mathop{\mathcal{H}\!\mathit{om}}\nolimits}
\def\SheafExt{\mathop{\mathcal{E}\!\mathit{xt}}\nolimits}
\def\Sch{\mathit{Sch}}
\def\Mor{\mathop{\mathrm{Mor}}\nolimits}
\def\Ob{\mathop{\mathrm{Ob}}\nolimits}
\def\Sh{\mathop{\mathit{Sh}}\nolimits}
\def\NL{\mathop{N\!L}\nolimits}
\def\CH{\mathop{\mathrm{CH}}\nolimits}
\def\proetale{{pro\text{-}\acute{e}tale}}
\def\etale{{\acute{e}tale}}
\def\QCoh{\mathit{QCoh}}
\def\Ker{\mathop{\mathrm{Ker}}}
\def\Im{\mathop{\mathrm{Im}}}
\def\Coker{\mathop{\mathrm{Coker}}}
\def\Coim{\mathop{\mathrm{Coim}}}

% Boxtimes
%
\DeclareMathSymbol{\boxtimes}{\mathbin}{AMSa}{"02}

%
% Macros for moduli stacks/spaces
%
\def\QCohstack{\mathcal{QC}\!\mathit{oh}}
\def\Cohstack{\mathcal{C}\!\mathit{oh}}
\def\Spacesstack{\mathcal{S}\!\mathit{paces}}
\def\Quotfunctor{\mathrm{Quot}}
\def\Hilbfunctor{\mathrm{Hilb}}
\def\Curvesstack{\mathcal{C}\!\mathit{urves}}
\def\Polarizedstack{\mathcal{P}\!\mathit{olarized}}
\def\Complexesstack{\mathcal{C}\!\mathit{omplexes}}
% \Pic is the operator that assigns to X its picard group, usage \Pic(X)
% \Picardstack_{X/B} denotes the Picard stack of X over B
% \Picardfunctor_{X/B} denotes the Picard functor of X over B
\def\Pic{\mathop{\mathrm{Pic}}\nolimits}
\def\Picardstack{\mathcal{P}\!\mathit{ic}}
\def\Picardfunctor{\mathrm{Pic}}
\def\Deformationcategory{\mathcal{D}\!\mathit{ef}}

% --- appended 2026-10-12 (Chief): the project own definitions of the reference/slogan/history
% display environments, copied mechanically from the shared Stacks preamble lines 8-12
% (\usepackage{verbatim} plus the three \newenvironment lines, which the chapter preamble keeps
% verbatim from the project preamble). Nothing added, removed or redefined: these environments
% carry no mathematical content and the project itself typesets them as comments. Without these
% three lines a counted statement whose own text uses them fails with Environment slogan undefined.
\usepackage{verbatim}
\newenvironment{reference}{\comment}{\endcomment}
\newenvironment{slogan}{\comment}{\endcomment}
\newenvironment{history}{\comment}{\endcomment}
\title{Power torsion modules become invariant under base change exactly when the quotients by the ideal agree}
\author{The Stacks Project Authors}
\date{stacks-project:more-algebra.tex:local-snapshot-20261012; GNU FDL 1.2 or later}
\begin{document}
\maketitle
\setcounter{subsection}{1}
\noindent\textit{Editorial note.} All delivered mathematical text is the author's own, copied from the pinned Stacks Project chapter file more-algebra.tex; no statement or proof step was rewritten, reordered or removed. The counted statement is the labeled lemma environment at source lines 24263-24273 (canonical tag 05E9) and the counted proof is the proof environment at source lines 24275-24298, which immediately follows it and is its own proof.\par\medskip
\noindent\textit{Editorial note.} The counted statement begins by assuming that the pair satisfies the equivalent conditions of the preceding lemma; that lemma's own statement is delivered as the context span of this card, so the hypothesis can be read and its label resolves inside the compiled document without a navigation stub.\par\medskip
\noindent\textit{Editorial note.} The preceding lemma's statement lies between two other items of the section, and the proof that follows it in the source is not part of this card; only its statement is delivered as context, which is why the context span ends at the end of that environment.\par\medskip
\noindent\textit{Editorial note.} The preamble span runs to line 90, the current end of the retained preamble, which includes the appended environment definitions; the counted text does not use them, and the span is taken whole so that the file version is unambiguous.\par\medskip
\noindent\textit{Editorial note.} The source URL is the source-level Stacks Project URL, as required when more than one card is produced from one source; the per-claim tag is carried by claim\_key and by this note. The counted statement and proof cite no bibliography entry, so this card needs no bibliography.\par\medskip
\noindent\textit{Allowed standard prerequisites (editorial).}\begin{itemize}
\item A module annihilated by the ideal is a module over the quotient, and tensoring it over the quotient with the base changed quotient gives back the module when the quotients agree. used at proof lines 24279-24284. Proves the base case of the induction, which is the heart of the equivalence.
\item Flatness of the ring map makes the tensor product with it exact, so a short exact sequence of modules gives a short exact sequence after tensoring. used at proof lines 24288-24293. Carries the induction step from the submodule and the quotient to the whole module.
\item A power torsion module is the union of the submodules annihilated by the successive powers of the ideal, and tensor products commute with colimits. used at proof lines 24294-24297. Passes from the modules annihilated by a fixed power to an arbitrary power torsion module.
\end{itemize}\medskip
\begin{lemma}
\label{lemma-characterize-flatness-on-torsion}
Let $\varphi : R \to S$ be a ring map. Let $I \subset R$ be an ideal.
The following are equivalent
\begin{enumerate}
\item $\varphi$ is flat and $R/I \to S/IS$ is faithfully flat,
\item $\varphi$ is flat, and the map
$\Spec(S/IS) \to \Spec(R/I)$ is surjective.
\item $\varphi$ is flat, and the base change functor
$M \mapsto M \otimes_R S$ is faithful on modules annihilated by $I$, and
\item $\varphi$ is flat, and the base change functor
$M \mapsto M \otimes_R S$ is faithful on $I$-power torsion modules.
\end{enumerate}
\end{lemma}

\begin{lemma}
\label{lemma-neighbourhood-isomorphism}
Assume $(\varphi : R \to S, I)$ satisfies the equivalent conditions of
Lemma \ref{lemma-characterize-flatness-on-torsion}.
The following are equivalent
\begin{enumerate}
\item for any $I$-power torsion module $M$, the natural map
$M \to M \otimes_R S$ is an isomorphism, and
\item $R/I \to S/IS$ is an isomorphism.
\end{enumerate}
\end{lemma}

\begin{proof}
The implication (1) $\Rightarrow$ (2) is immediate.
Assume (2). First assume that $M$ is annihilated by $I$.
In this case, $M$ is an $R/I$-module. Hence, we have an isomorphism
$$
M \otimes_R S = M \otimes_{R/I} S/IS = M \otimes_{R/I} R/I = M
$$
proving the claim. Next we prove by induction that $M \to M \otimes_R S$
is an isomorphism for any module $M$ is annihilated by $I^n$. Assume
the induction hypothesis holds for $n$ and assume $M$ is annihilated by
$I^{n + 1}$. Then we have a short exact sequence
$$
0 \to I^nM \to M \to M/I^nM \to 0
$$
and as $R \to S$ is flat this gives rise to a short exact sequence
$$
0 \to I^nM \otimes_R S \to M \otimes_R S \to M/I^nM \otimes_R S \to 0
$$
Using that the canonical map is an isomorphism for $M' = I^nM$ and
$M'' = M/I^nM$ (by induction hypothesis) we conclude the same thing is
true for $M$. Finally, suppose that $M$ is a general $I$-power torsion
module. Then $M = \bigcup M_n$ where $M_n$ is annihilated by $I^n$
and we conclude using that tensor products commute with colimits.
\end{proof}

\end{document}
